% Item 7(a) do parecer — trechos para integrar à dissertação.
% Fonte única TFLM: coleta 20261009T120557541292Z, 2.000 imagens.
% Não compilar este arquivo isoladamente: são alternativas/trechos de integração.

% 1. Em tab:tempos-inferencia-aceleracao (Tabela 5.2), SUBSTITUIR a linha
% TFLM aproximada já existente. Não adicionar uma segunda linha para o runtime.
\shortstack[l]{TFLM + ESP-NN,\\ESP32-CAM, 240 MHz} &
1,777 s &
1,761 s &
1,793 s &
81,06$\times$ \\

% 2. SUBSTITUIR a nota atual da Tabela 5.2, que diz não haver mínimo/máximo.
\begin{minipage}{\textwidth}
\small \textbf{Nota:} Os valores de TFLM com ESP-NN foram calculados sobre
2.000 imagens da coleta \texttt{20261009T120557541292Z}. O tempo inclui a
preparação da entrada e a execução de \texttt{Invoke}, excluindo download,
gravação de checkpoints e leitura da saída. A aceleração de 81,06$\times$
utiliza a referência interpretada a 160~MHz da tabela e é uma razão
descritiva entre cenários com frequências e tamanhos de amostra diferentes.
A comparação de latência entre WASM AOT e TFLM a 240~MHz não isola o custo
exclusivo da portabilidade, pois os kernels e as otimizações diferem.
\end{minipage}

% Alternativa: se a coluna de aceleração for reservada às estratégias WAMR,
% usar --- no lugar de 81,06$\times$ e explicar essa delimitação na nota.

% 3. Em tab:acuracia-configuracoes (Tabela 5.3), SUBSTITUIR a linha
% "TFLite ESP32-CAM". Os valores numéricos permanecem iguais.
\shortstack[l]{TFLM + ESP-NN,\\ESP32-CAM, 240 MHz} &
1.990 & 9 & 1 & 99,50\% \\

% 4. Acrescentar junto à Tabela 5.3.
Os resultados de TFLM com ESP-NN nas
Tabelas~\ref{tab:tempos-inferencia-aceleracao}
e~\ref{tab:acuracia-configuracoes} utilizam a mesma coleta,
identificada por \texttt{20261009T120557541292Z}.
Foram processadas 2.000 imagens sem falhas de processamento, com
1.990 acertos, nove predições válidas incorretas e um empate.
O empate foi contabilizado como ocorrência sem predição válida e
como não correto no cálculo da acurácia, cujo denominador é 2.000.

% 5. SUBSTITUIR o parágrafo comparativo que usa apenas "aproximadamente 1,7".
A comparação de desempenho com uma referência embarcada utilizou
o mesmo modelo de sonolência no ESP32-CAM a 240~MHz, executado com
TensorFlow Lite Micro, componente versão 1.3.5, e ESP-NN habilitado.
O tempo médio foi de 1,776715~s por imagem, com mínimo de 1,760819~s
e máximo de 1,792774~s. O intervalo compreende a preparação da entrada
e a execução de \texttt{Invoke}, excluindo download, gravação de
checkpoints e leitura da saída. Não houve aquecimento explícito,
conforme o metadata da coleta.
Em comparação com os 6,195~s do WASM AOT a 240~MHz apresentados na
Tabela~\ref{tab:tempos-inferencia-aceleracao}, a latência média do
TFLM foi aproximadamente 3,49 vezes menor. Essa razão caracteriza
o desempenho relativo das implementações avaliadas, sem isolar
um custo exclusivo de portabilidade, pois diferem os kernels,
as otimizações e a organização da execução. A comparação entre
interpretação e AOT do WAMR a 160~MHz permanece a avaliação específica
do efeito da estratégia de execução WebAssembly no mesmo hardware
e frequência.

% 6. Opcional: acrescentar também a configuração à Tabela 5.1
% (tab:configuracoes-experimentos-wasm), cujo formato possui seis colunas.
TFLM + ESP-NN &
ESP32-CAM / TFLM 1.3.5 &
240 MHz &
4 MiB PSRAM &
2.000 &
1,777 s \\
